\paragraph{Onde o \texttt{volatile} aparece no bytecode?}
A classe abaixo tem um campo \texttt{volatile} usado como flag, um campo
\texttt{final} e duas formas de inicialização preguiçosa: um DCL com o
campo \texttt{instancia} declarado \texttt{volatile} e o idiom do
holder, com a classe aninhada \texttt{Portador}. A listagem usa
\texttt{javap -v -p}, que mostra as flags de cada campo.

\lstinputlisting[language=Java]{generated/source/PublicacaoSegura.java}

\inputbytecode{PublicacaoSegura}

\begin{itemize}
\item As flags dos campos estão logo depois do constant pool.
  \texttt{instancia} tem \texttt{(0x004a) ACC\_PRIVATE, ACC\_STATIC,
  ACC\_VOLATILE}, \texttt{limite} tem \texttt{(0x0012) ACC\_PRIVATE,
  ACC\_FINAL} e \texttt{ativo} tem \texttt{(0x0042) ACC\_PRIVATE,
  ACC\_VOLATILE}. O valor é a soma das flags da Tabela 4.5-A da JVMS
  SE~25 (§4.5): \texttt{ACC\_PRIVATE} vale \texttt{0x0002},
  \texttt{ACC\_STATIC} \texttt{0x0008}, \texttt{ACC\_FINAL}
  \texttt{0x0010} e \texttt{ACC\_VOLATILE} \texttt{0x0040}. A tabela
  descreve \texttt{ACC\_VOLATILE} como \enquote{Declared volatile;
  cannot be cached} e \texttt{ACC\_FINAL} como \enquote{never directly
  assigned to after object construction (JLS §17.5)}. Um campo não pode
  ter as duas flags.
\item Em \texttt{obter}, a primeira leitura de \texttt{instancia} é o
  \texttt{getstatic} do offset~0, e a escrita é o \texttt{putstatic} do
  offset~32. O \texttt{getstatic} é o mesmo que
  \texttt{obterPorPortador} usa para ler \texttt{Portador.INSTANCIA},
  que não é \texttt{volatile}. No construtor, os \texttt{putfield} dos
  offsets 6 e 11 gravam \texttt{ativo}, que é \texttt{volatile}, e
  \texttt{limite}, que é \texttt{final}, com a mesma instrução. O que
  marca o campo é a flag na declaração, não a instrução.
\item O \texttt{ifnonnull} do offset~5 pula para o offset~45 quando a
  primeira leitura já vê o objeto: é o caminho sem lock. Senão, o
  \texttt{monitorenter} do offset~12 trava o objeto \texttt{Class}, o
  \texttt{getstatic} do offset~13 lê o campo de novo e o
  \texttt{ifnonnull} do offset~18 pula a criação se outra thread já a
  fez. O \texttt{monitorexit} está nos offsets 36 e 42; a tabela de
  exceções desvia os offsets 13--37 para o 40, que libera o monitor e
  relança a exceção.
\item O objeto é criado nos offsets 21--27, e o \texttt{invokespecial}
  do offset~27 executa o construtor. Na thread que cria, as gravações
  do construtor vêm antes do \texttt{putstatic} do offset~32 na ordem
  do programa. A §17.4.5 da JLS SE~25 dá happens-before entre ações da
  mesma thread em ordem de programa e entre a escrita \texttt{volatile}
  e as leituras subsequentes do campo, e a relação é transitiva. O FAQ
  da JSR~133 explica o conserto do DCL por essa via: com o campo
  \texttt{volatile}, há happens-before entre a inicialização do objeto
  pela thread que o constrói e a devolução do valor pela thread que o
  lê.
\item \texttt{obterPorPortador} é um único \texttt{getstatic} de
  \texttt{Portador.INSTANCIA}, sem lock no código. A classe aninhada
  aparece nos atributos \texttt{NestMembers} e \texttt{InnerClasses}, no
  fim da listagem. O FAQ da JSR~133 recomenda esse idiom e diz que ele
  é correto pelas garantias de inicialização de campos \texttt{static}:
  um campo atribuído num inicializador \texttt{static} fica visível,
  corretamente, a toda thread que acessa a classe.
\end{itemize}

\paragraph{O que isso muda no seu código.}
Cada recomendação tem seu escopo.
\begin{itemize}
\item \emph{Java~5 (2004) em diante:} use \texttt{volatile} para uma
  flag que uma thread escreve e outras leem, como \texttt{ativo} acima,
  e para publicar uma referência. A escrita \texttt{volatile}
  happens-before as leituras subsequentes (JLS SE~25, §17.4.5). No
  modelo anterior, \texttt{volatile} podia ser reordenado com acessos
  comuns, o que minava esse uso, segundo o FAQ. \texttt{volatile} não
  torna atômico um ler-somar-gravar como \texttt{contador++}
  (seção~\ref{sec:modelo-memoria}).
\item \emph{Java~5 (2004) em diante:} para um objeto imutável
  publicado entre threads, declare os campos \texttt{final}. O FAQ diz
  que os valores atribuídos a campos \texttt{final} no construtor ficam
  visíveis a todas as threads sem sincronização, desde que o objeto
  seja construído corretamente, isto é, sem deixar a referência a ele
  escapar durante a construção: nada de atribuí-la a um campo
  \texttt{static} ou registrá-la como listener dentro do construtor. A
  garantia está na §17.5 da JLS SE~25. No modelo antigo, campos
  \texttt{final} não tinham essa garantia (FAQ).
\item \emph{Java~5 (2004) em diante:} para um singleton
  \texttt{static} preguiçoso, prefira o holder; o FAQ o descreve como
  \enquote{thread-safe and a lot easier to understand}. Se usar DCL, o
  campo tem de ser \texttt{volatile}. Segundo a Declaração, isso
  \enquote{does not work under JDK4 and earlier}; nessas versões, nem
  o \texttt{volatile} conserta o idiom.
\end{itemize}

\paragraph{Onde aprofundar.}
O modelo atual, com o exemplo do \texttt{contador++}, está na
seção~\ref{sec:modelo-memoria}; \texttt{monitorenter} e
\texttt{monitorexit}, na seção~\ref{sec:sincronizacao}. O FAQ, a
especificação e o cookbook da JSR~133 estão na página de Pugh, o site
oficial da JSR (\citetitle{doc:pugh-jmm-pagina}). Para o modelo
formal, veja \citetitle{paper:java-memory-model}; para as consequências
práticas, \citetitle{artigo:shipilev-jmm} e \citetitle{livro:goetz2006}.
As classes de \texttt{java.util.concurrent} são o tema da
seção~\ref{sec:juc}.
